\section{Limitations and Future Work}\label{sec:limitations}

A promising direction for future work is extending SCONE beyond short $n$-grams by enabling caching for longer queries. A central challenge lies in designing effective keys for such queries. Using raw text as keys may result in low cache hit rates, as semantically similar queries often vary at the surface level. Alternatively, using semantic embeddings as keys would require discretization techniques to map continuous embeddings into a discrete key space that supports efficient indexing.

One limitation of the current study is that we evaluate SCONE only on models with up to 3B parameters at training time. This constraint is primarily due to hardware limitations, which restrict our ability to conduct large-scale pretraining on larger models. Although the model scales we tested are widely used in real-world applications, exploring the performance of SCONE at larger scales presents an exciting direction for future research. We believe applying SCONE to larger models would be highly beneficial to the community and leave this exploration for future work. Notably, our experiments in \Cref{sec:exps_openwebtext} provide encouraging evidence, as the benefits of SCONE are consistent across all model sizes we tested.

\section{Additional Related Work}\label{sec:add_related_work}

\paragraph{Contextualized Word Embeddings.} Words can have different meanings depending on context. Prior work has incorporated context into word embeddings, either from the entire sequence \citep{mccann2017learned,peters2018deep} or short \ngram{n}s \citep{gupta2019better}, before applying them to downstream tasks. Modern language models inherently use contextualized token embeddings, leveraging attention mechanisms. In this study, we extend the embedding layer to include contextualized f-gram embeddings for each token. A key novelty is that our approach allows embeddings to be precomputed and offloaded from accelerators, providing contextual embeddings for each token without increasing FLOPS and accelerator memory usage at inference.


\paragraph{Tokenization in Language Models.} Our method assumes a predefined vocabulary from a trained tokenizer. Several popular algorithms exist for training tokenizers \citep{sennrich2015neural, wu2016google, kudo2018subword, kudo2018sentencepiece}. In this work, we use a BPE tokenizer, following prior seminal works \citep{radford2019language, touvron2023llama}. However, our method is not tied to any specific tokenization algorithm and can be applied seamlessly to others.

Tokenization-free language models have also been widely explored \citep{kim2016character, choe2019bridging, xue2022byt5, yu2023megabyte, wang2024mambabyte, deiseroth2024t, meta2024large, pagnoni2024byte}. While we have not tested our method on tokenization-free models, we believe our core idea—introducing an off-accelerator embedding layer by precomputing embeddings for frequent input patterns—remains applicable.



\paragraph{Mixture-of-Experts (MoE) and Memory Layers.} MoE and memory layers are two established approaches for scaling language models within a fixed FLOPS budget.

MoE layers replace traditional feedforward layers with multiple parallel ``experts," activating only one or a few per token using a lightweight routing mechanism \citep{shazeer2017outrageously,lepikhin2021gshard,fedus2022switch,jiang2024mixtral,he2024mixture}. This allows the model to scale by increasing the number of experts without increasing the computational cost per token. However, all experts must reside on the accelerator, resulting in significantly higher memory usage.

Memory layers, on the other hand, store large collections of embeddings (continuous vectors) and retrieve the nearest neighbors during the forward pass via (approximate) similarity search \citep{weston2014memory,sukhbaatar2015end,lample2019large,berges2024memory}. These retrieved embeddings enhance the model’s capabilities without adding much to the FLOPS budget. Despite improvements in similarity search techniques \citep{lample2019large,johnson2019billion}, memory layers still require storing the embeddings on the accelerator, making memory demands impractically high at larger scales \citep{berges2024memory}. Furthermore, because embeddings are typically updated via backpropagation, memory layers introduce additional challenges related to sparse updates as the memory size grows.

While MoE layers, memory layers, and our proposed method \SCONE all maintain fixed inference FLOPS, a key advantage of \SCONE is it also maintains fixed accelerator memory usage at inference. By focusing on the input embedding layer, \SCONE ensures that the computational overhead remains $O(1)$ during inference and enables offloading the additional parameters off-accelerator with negligible impact on end-to-end latency.



\paragraph{Implicit $n$-gram Patterns in Transformers.} Recent work analyzing the internal mechanisms of transformers has shown that these models often utilize implicit $n$-gram patterns for prediction \citep{geva2020transformer,geva2022transformer,voita2023neurons}. For instance, \citet{chen2024jet} demonstrate that certain attention heads can detect specific $n$-gram patterns, while MLPs can perform linguistic operations such as adding the “-ing” suffix. These findings underscore the importance of $n$-gram information in language modeling and offer a potential explanation for the effectiveness of \SCONE. An interesting future direction is to examine how \SCONE's f-gram embeddings interact with the transformer’s implicit $n$-gram patterns.

\paragraph{Embedding Sparsity in Multilingual Applications and Recommender Systems.} This work focuses on a common setting for training LLMs: language modeling on large-scale text corpora, primarily in English. However, scaling embedding layers presents challenges beyond this context, particularly due to frequency-related performance degradation caused by sparsity. Multilingual applications are one such scenario. Two phrases in different languages may refer to the same concept but correspond to different embedding vectors. Their embeddings should ideally be close. Recent work has explored methods for learning transferable embeddings in cross-lingual settings \citep{artetxe2019cross, chen2023improving}. Another relevant example is scaling the embeddings for recommender systems \citep{chen2019lambdaopt, liu2021learnable}, where embeddings often dominate the model's parameter count due to the high cardinality of user or item categories. For both scenarios, \SCONE’s strategy, i.e., parameterizing large embedding tables using a neural network, provides a complementary approach to help mitigate sparsity issues.